\section{Limitations}
\label{sec:limitations}

The findings of this paper are subject to several methodological limitations
that constrain the strength of the conclusions drawn. This section identifies
the most consequential threats to validity, explains their implications for the
central arguments, and suggests how future work might address them.

% ============================================================
\subsection{Correlation and Causation}
\label{sec:lim-causation}
% ============================================================

The analytical framework employed throughout this paper documents statistical
associations between offensive and defensive quality metrics and competitive
outcomes. It does not establish causal relationships. The finding that offensive
\xvoa{} explains more variance in winning than defensive \xvoa{} demonstrates
that offensive execution covaries more strongly with success, but it cannot prove
that investing in offensive talent \emph{causes} teams to win more games. The
distinction matters because resource allocation decisions -- the strategic
implications discussed in \Cref{sec:implications} -- require causal reasoning,
yet the evidence base is purely observational.

A specific concern is reverse causation operating through game state. Teams that
establish early leads often shift to more conservative play-calling: running the
ball to consume clock, avoiding aggressive downfield passing, and accepting
shorter drives that reduce turnover risk. This behavioral shift means that
teams already winning may generate different \xscore{} profiles not because
their offensive quality is superior, but because their game situation permits
lower-variance strategy. The \gds{} framework mitigates this partially by
measuring per-drive execution quality rather than volume -- a team running
conservative drives still has each drive evaluated against the expected scoring
probability from its starting field position. However, this mitigation is
incomplete. If conservative play-calling systematically produces drives that
exceed expectation (for example, by avoiding turnovers that the model expects at
league-average rates), then \xvoa{} values for leading teams may be inflated.
A specific mediation concern -- that defensive quality generates favorable
starting field position which inflates offensive \xvoa{} -- was directly tested
in \Cref{sec:mediation-findings}. The mediation analysis found that only 9.2\%
of the offensive \xvoa{}-to-wins association is attributable to
defense-generated field position, with 90.8\% of the offensive value operating
independently. This result substantially mitigates the indirect-pathway
version of the causal concern, though it does not resolve the broader issue
of reverse causation through game state.

A randomized experiment -- randomly assigning offensive and defensive talent to
teams -- is impossible in professional sport. Observational data can establish
the pattern of association, as this paper has done, but the causal interpretation
remains an inference rather than a demonstration.

% ============================================================
\subsection{Small Championship Sample}
\label{sec:lim-sample}
% ============================================================

The logistic regression analysis presented in \Cref{sec:logistic-results}
attempts to model the probability of winning the Super Bowl as a function of
offensive and defensive \xvoa{}. The fundamental constraint is that the
2018--2024 study window contains exactly seven Super Bowl winners across seven
seasons. With only seven positive events and two predictors, the
events-per-variable (EPV) ratio is 3.5, well below the standard recommendation
of 10 EPV for logistic regression \parencite{peduzzi1996simulation}. This low EPV
means that coefficient estimates are unstable, confidence intervals are wide,
and the model has limited statistical power to detect effects of moderate size.

Although both the offensive ($p = 0.003$) and defensive ($p = 0.027$)
coefficients achieve statistical significance in the current model, the low EPV
means that these estimates carry substantial uncertainty and should not be
over-interpreted in terms of their relative magnitudes. The wider standard error
on the defensive coefficient (SE $= 0.178$ vs $0.126$ for offense) reflects the
inherent imprecision of a seven-event logistic model. This limitation is
structural rather than methodological -- no analytical choice can increase the
number of Super Bowl winners per season. The pre-registered hypothesis tests
and descriptive analyses that form the primary evidential basis of this paper
do not suffer from this constraint to the same degree, as they operate on the
full 94 playoff team-seasons rather than only the 7 champions. Nevertheless,
any claim about championship-specific effects must be understood as provisional
until more seasons of data become available.

% ============================================================
\subsection{Schedule Strength}
\label{sec:lim-opponent}
% ============================================================

The \gds{} framework measures raw execution quality: how much value a team's
offense or defense generates relative to league-average expected outcomes, without
adjusting for the quality of opposition faced at the drive level. A defense
achieving high \xvoa{} against a schedule of below-average offenses is rated
identically to one achieving the same \xvoa{} against elite offenses. This
absence of drive-level opponent adjustment introduces potential bias if schedule
strength correlates systematically with team archetype.

To assess whether this limitation materially affects the archetype findings, I
constructed a leave-one-out opponent-strength variable and included it as a
control in the OLS regression predicting playoff wins
(\Cref{sec:statistical-confirmation}). The opponent-strength coefficient was
effectively zero ($-0.001$, $p = 0.985$), and neither the offensive nor
defensive \xvoa{} coefficients changed meaningfully when schedule strength was
controlled. This null result provides empirical reassurance that the documented
offensive advantage is not an artifact of schedule imbalance.

Nevertheless, a more principled drive-level opponent adjustment -- in which each
drive's expected baseline accounts for the opposing unit's quality -- represents
a natural extension. Such an adjustment would require an iterative estimation
procedure (since each team's quality depends on the quality of teams it has
faced, which in turn depends on the quality of teams \emph{they} have faced),
analogous to the iterative strength-of-schedule corrections used in Elo rating
systems. This extension was beyond the scope of the present paper but would
further strengthen the internal validity of the \gds{} decomposition at the
drive level, even though the season-level robustness check suggests the bias is
negligible.

% ============================================================
\subsection{Data Limitations}
\label{sec:lim-data}
% ============================================================

The \texttt{nflfastR} play-by-play dataset provides comprehensive
situational information (down, distance, field position, score
differential, time remaining) but lacks several categories of data that
would enrich the analysis.
Player tracking data (Next Gen Stats), which records the position and velocity
of all 22 players at 10~Hz, is not publicly available for research use. Without
tracking data, the model cannot distinguish between value generated by scheme
design and value generated by individual talent execution. A team might achieve
high offensive \xvoa{} because its offensive coordinator designs plays that
create favorable situations, or because its quarterback makes exceptional
decisions under pressure, or both -- the current framework cannot disaggregate
these contributions.

Additionally, personnel groupings (e.g., 11~personnel, 12~personnel), pre-snap
alignment information, and defensive coverage classifications are absent from
the modeling features. Weather conditions, player injuries, and weekly roster
availability -- all of which influence drive outcomes -- are not incorporated. The
practical effect is that \xscore{} predictions reflect league-average
relationships between situational variables and scoring probability, without
conditioning on the specific personnel or environmental context of each drive.
This omission likely adds noise to drive-level predictions but should not
introduce systematic bias at the season level, provided that personnel and
weather effects are approximately symmetrically distributed across team
archetypes.

% ============================================================
\subsection{Temporal Confound: The Rules Environment}
\label{sec:lim-temporal}
% ============================================================

The 2018--2024 study window coincides with an NFL rules environment that
explicitly favors offensive production. The 2018 season introduced expanded
roughing-the-passer protections \parencite{pereira2018roughing}, and subsequent seasons have seen increased
enforcement of defensive holding, illegal contact, and player safety rules that
limit the physicality available to defensive players. These rule changes have
contributed to historically high passing efficiency, scoring rates, and offensive
production across the league during the study period. The observed offensive
advantage in predicting success may therefore be partially an artifact of the
regulatory environment rather than a fundamental structural truth about the
sport of American football.

If future rule changes were to rebalance the equilibrium between offense
and defense -- for example, by permitting more physical coverage techniques
or reducing protections for quarterbacks outside the pocket -- the magnitude
of the offensive advantage documented here might attenuate. The
out-of-distribution validation conducted on 2014--2017 data
(\Cref{sec:data}) demonstrates that the \xscore{} model generalizes well
to a moderately different rules environment, providing confidence in the
model's technical validity. However, the archetype analysis
and hypothesis tests were not extended to pre-2018 seasons, meaning that the
question of whether offense has \emph{always} been more important or whether
this is a modern phenomenon remains unanswered. A longitudinal extension
spanning multiple regulatory eras (for example, 2002--2024) would be required to
disentangle the structural contribution of offense from the rule-induced
amplification of offensive metrics.

% ============================================================
\subsection{Special Teams as a Residual Quantity}
\label{sec:lim-special-teams}
% ============================================================

The \gds{} framework decomposes team quality into three components:
offensive \xvoa{}, defensive \xvoa{}, and a special teams value term
(ST\_Value). As described in \Cref{sec:gds}, the special teams component
measures field-position deviations from transition-type baselines: for each
drive, ST\_Value equals the starting \xscore{} xEP minus the league-average xEP
for drives of the same transition type (e.g., post-kickoff, post-punt). This
construction captures genuine special teams impact on field position (kick
returns, punt coverage) but does not independently model special teams scoring
events. It also conflates any systematic bias in the \xscore{} baseline
predictions with special teams value, since both enter through the starting
xEP term. In consequence, the magnitude and
variability of ST\_Value cannot be interpreted purely as a measure of special
teams quality. A dedicated special teams probability model -- estimating expected
points from kick-off situations, punt returns, and field goal attempts
independently -- would provide a more principled decomposition. Such a model was
beyond the scope of this paper, and the residual approach was judged acceptable
given that special teams contribute a relatively small share of total game value
in the modern NFL. Nevertheless, any future extension of the \gds{} framework
should prioritize replacing the residual with an independently estimated special
teams component.
